\section{Scope, questions, and evidence boundary}

\subsection{Scientific question}

The study asks a representation-sensitive question inside one fixed rank-two physical
subspace.  Suppose two matrix fields are related pointwise by a unitary change of
retained basis and therefore have identical ordered eigenvalues at every crystal
momentum.  Must they have the same locality when each is Fourier transformed in its own
declared retained frame?  The answer is no: pointwise spectral equivalence constrains
eigenvalues but not the momentum dependence of the representing basis.  A smooth,
momentum-dependent gauge action can preserve every band energy while redistributing
matrix-valued hopping power over cell displacements.

The numerical study makes that statement quantitative in four separate channels:
\begin{enumerate}
 \item a bounded parent-hopping family, which changes the ambient parent dispersion
       while preserving all other controls;
 \item nested momentum-mesh refinement, which measures finite-mesh stability and
       residual consistency without changing the physical model;
 \item deterministic adaptive amplitude refinement, which measures interpolation
       defects in the gauge-attack parameter and requires disjoint validation before
       convergence; and
 \item ordinary-lattice truncation, which measures omitted hopping power, represented
       Hamiltonian error, and sorted-band error on a mesh disjoint from the construction
       mesh.
\end{enumerate}
These channels are retained together but are not collapsed into one error number.
Parent-model variation, reciprocal discretization, parameter interpolation, and model
reduction answer different questions.

\subsection{What ``band equivalent'' means here}

For the reference retained Hamiltonian $H(q)$ and comparison Hamiltonian $A_a(q)$,
band equivalence means
\begin{equation}
 \operatorname{spec} A_a(q)=\operatorname{spec} H(q)
 \quad\text{for every evaluated }q,
 \label{eq:band-equivalence-definition}
\end{equation}
with eigenvalues compared in sorted order.  It does not mean $A_a(q)=H(q)$ as matrices,
it does not identify their Fourier coefficients, and it does not permit an arbitrary
momentum-dependent alignment during the comparison.  The scientific comparison is
intentionally restricted to one momentum-independent real rotation, with a separately
optimized nonnegative scalar where that scalar is part of the declared question.

This restriction is substantive.  If the same momentum-dependent gauge transformation
used to construct $A_a$ were admitted during alignment, the comparison would be
trivial: the fields could be returned to pointwise equality by construction.  The
restricted problem instead asks what can be represented by one global rank-two
coordinate choice.

\subsection{Finite-control status}

Everything computed in this study is finite.  The parent is a synthetic $4\times4$
Hermitian matrix field.  Its retained subspace has rank two.  Construction uses a finite
centered uniform momentum mesh.  Fourier coefficients form a finite cyclic inventory
before an explicit even-mesh cyclic-to-ordinary conversion.  Amplitude studies use a
finite interval $0\le a\le\pi/4$, and truncation studies use finitely many integer
cutoffs.  No thermodynamic limit, continuum limit, complete extremum certificate, or
material calibration follows merely from these finite operations.

Two unsealed execution directories are presently retained.  The first preserves a
failed independent-verification attempt and has no figure inventory.  The authorized
no-overwrite successor preserves the same canonical calculation bytes, passed aggregate
independent verification at the unchanged installed tolerance, and produced the closed
six-figure inventory through stored admission.  Neither directory is a reviewed release,
a sealed evidence package, a publication result, or a material prediction.  Section~5
keeps the failed verification history, the successful bounded route agreement, and the
remaining unresolved scientific channels distinct.

\subsection{Integrity, verification, and scientific validity}

Three evidentiary notions are kept distinct:
\begin{description}
 \item[Canonical integrity.] Canonical typed JSON and SHA-256 identify exact bytes and
 reject alteration or noncanonical re-encoding.  A digest says nothing by itself about
 whether a formula is physically appropriate.
 \item[Independent route agreement.] A verifier reconstructs the same finite observables
 through direct formulas that avoid the production scientific owners.  Passing would
 establish bounded numerical agreement under the installed comparison policy.
 \item[Scientific validity.] Applicability of the model, adequacy of the controls,
 convergence of relevant limits, interpretation of observables, and suitability for a
 scientific claim require arguments beyond byte identity and software-route agreement.
\end{description}
The first candidate satisfies neither a passing aggregate-verification claim nor a
rendered-artifact claim; its failure remains a diagnostic fact.  The successor satisfies
the bounded independent-route and authenticated-rendering contracts, but that software
agreement does not upgrade the finite calculation to scientific validation or sealed
evidence.

\subsection{Ownership boundary}

Quantum ESPRESSO and Wannier90 are not used by this synthetic study.  The study does not
replace either package.  Project Koios Optimization supplies unit-free numerical
algorithms, PhysKit supplies physical and finite-periodic semantics, and
\texttt{ksdft2effmass} owns the installed model, policy thresholds, closed outcome
variants, aggregate retention, verification admission, and calculation-facing
visualization semantics.  This division matters because independent verification must
not call the same production owner whose result it is meant to check.
